\section{Organisation}
\subsection{Übung}
\begin{itemize}
    \item \fatsf{Ziel:} Wiederholung/Vertiefung des Vorlesungsmaterials
    \item \fatsf{Inhalt der Übungen:}
    \begin{itemize}
        \item Übungsaufaben ähnlich zur Klausuraufgaben $\rightarrow$ ideale Vorbereitung auf die Klausur
        \item Themen sind \fatsf{Klausurrelevent}
    \end{itemize}
    \item Zweiwöchentliche Übungsblätter
    \item Eine Abgabe und Korrektur durch uns ist nicht vorgesehen
    \item Übungen: vor Ort und via Zoom
    \item Erste Übungen: 28.10.2024 - 01.11.2024
    \item Sprechstunden vor Ort und via Zoom altermierend mit Übungen
\end{itemize}
\subsection{Klausur}
\begin{itemize}
    \item Termin: \fatsf{11. März 2025, vorläufig!}
    \item Schriftliche Prüfung von 90 Minuten Dauer
    \item Wir planen rechtzeitig zum Lernen eine \fatsf{Probeklausur} herauszugeben
\end{itemize}
\subsection{Bonusübung}
\begin{itemize}
    \item \fatsf{Freiwillige} Hausübung über die Weihnachtspause
    \begin{itemize}
        \item Details und Termine folgen in Moodle
    \end{itemize}
    \item Es können 5 Bonuspunkte (ein Notenschritt) für die Klausur erreicht werden
    \item \fatsf{Gruppenarbeit ist nicht erlaubt}
\end{itemize}
\subsection{Terminübersicht}
\begin{itemize}
    \item Mittwoch 9:50-11:30
    \item Donnerstag 9:50-11:30
\end{itemize}
\subsection{Lernziele}
\begin{itemize}
    \item Sie wissen, \fatsf{Sicherheit ist fundamental} für jedes Computersystem
    \item Sie können \fatsf{Bedrohungen} und deren mögliche \fatsf{Auswirkungen} abschätzen
    \item Sie wissen wie \fatsf{Angriffe in der Praxis} funktionieren
    \item Sie können das Potential und die Limitierungen heutiger \fatsf{Sicherheitsmaßnahmen} und \fatsf{-technologien} einschätzen
    \item Sie kennen \fatsf{Gegenmaßnahmen} für ausgewählte Angriffe
    \item Sie wissen welche Fehler Sie \fatsf{nicht} machen sollten!
\end{itemize}
\subsection{Vorlesungsinhalte}
\begin{itemize}
    \item Einführung und Grundlagen
    \begin{itemize}
        \item Grundlegende Ideen, Begrifflichkeiten, Sicherheitsprinzipen und Eigenschaften
    \end{itemize}
    \item Kryptographie
    \begin{itemize}
        \item Verschlüsselung, Schlüsselaustausch, digitale Signaturen, ...
    \end{itemize}
    \item Netzwerksicherheit
    \begin{itemize}
        \item Angriffe und Gegenmaßnahmen im ISO/OSI-Schichtenmodell, TLS, Wifi, TOR, ...
    \end{itemize}
    \item Websicherheit
    \begin{itemize}
        \item Cookies, XSS, SQL-Injections, ...
    \end{itemize}
    \item Systemsicherheit
    \begin{itemize}
        \item Schadsoftware, Seitenkanalangriffe, Buffer Overflows, ...
    \end{itemize}
\end{itemize}